\documentclass[../main]{subfiles}
\begin{document}
\chapter{Interacción electromagnética}
\info{
Contenido}{
Interacción entre la corriente eléctrica y campo magnético, generación de f.e.m y fuerzas.
f.e.m inducida por un conductor que se mueve en un
campo magnético}
\info{Fórmulas}{
\begin{center}
 \begin{tabular}{c c}
 $\epsilon= B l v =f.e.m = N \frac{\Delta \Phi}{\Delta t}$ & $\epsilon_{auto}=\dfrac{\Delta \Phi}{\Delta t}$ \\
  & \\
 $e_{auto}=L\dfrac{\Delta I}{\Delta t}$ & $L= N \dfrac{\Phi}{I}$ \\
  & \\
 $F= B l I$ & $\Phi = B s$
 \end{tabular}
 \end{center}

siendo:

$s$: sección ($m^2$)

$\Phi$: Flujo magnético. (Tesla $m^2$)

$\Delta \Phi =\Phi_{final} -\Phi_{inicial} $: Variación de flujo magnético.

$\Delta t=t_{final}-t_{inicial}$:Variación de tiempo.

$I$: Corriente (Ampere)

$\Delta I = I_{final}-I_{inicial}$: Variación de corriente. (Ampere)

$L$: Coeficiente de autoinducción.

$N$: Número de espiras.

$F$: Fuerza (Newtons)

$\epsilon$: f.e.m. inducida.(Volts)

$\epsilon_{auto}$: f.e.m. de auto inducción.(Volts)


$B$: Inducción magnética.(Tesla)

$l$: Longitud del conductor.(metros)

$v$: Velocidad del conductor (m/s)

}
\iffalse
\begin{tcolorbox}
Teoría
\tcblower
\begin{tabular}{|p{12cm}|c|}\hline
Descripción & Código \\ \hline
 f.e.m inducida por un conductor  que se mueve en un campo magnético & \qrcode{https://www.youtube.com/watch?v=6N7ukJxvQ_M}\\\hline
Coeficiente de autoinducción & \qrcode{https://www.youtube.com/watch?v=Xb5GHlaWv7g}  \\\hline
Fuerza producida en un conductor por el que circula corriente dentro en un campo magnético &
\qrcode{https://www.youtube.com/watch?v=QvBbUOn_T4k}\\ \hline
\end{tabular}
 
\end{tcolorbox}
\fi

\actividad{Resolver los siguientes problemas}

\begin{enumerate}
    \item Un conductor se desplaza a una velocidad lineal de $v=10m/s$ dentro de un campo magnético fijo de $B=1,5T$ de inducción. Determinar el valor de la f.e.m. inducida en el mismo si posee una longitud de $l=1m$.
    \item A que velocidad se deberá mover un conductor de longitud $l=0.5m$ para que dentro de una inducción de campo magnético de $B=1T$ produzca una f.e.m de $e=5V$.
    \item Determinar la inducción de campo magnético $B$, que será necesaria para que un conductor de $l=1m$ de longitud que se mueve a $v=10m/s$ cortando las líneas del campo magnético produzca una f.e.m de $e=12V$.
    \item ¿Cuál la f.e.m de auto-inducción si el flujo magnético $\Phi$ que existe en el campo magnético producido por una bobina de una electroválvula si tiene esta un núcleo cilíndrico de diámetro $d=4cm$ y la inducción magnética en la misma crece de $B=0T$ a $B=1.5T$ en $\Delta t=0.3s$?
    \item Calcular el valor de la f.e.m de auto-inducción que se desarrollará una boina con coeficiente de auto-inducción de 50 milihernios si se aplica una corriente que crece desde cero hasta $I=20A$ en un tiempo de $\Delta t=0,1 seg.$
    \item Cuanto tiempo habrá que esperar para que la f.e.m de auto-inducción se reduzca a $e_{auto}=1V$ si se hace circular $I=10A$ por una bobina con un coeficiente de auto-inducción de $L=60 mH$  
    \item Cual es la f.e.m de auto inducción del problema anterior cuando pasa 1 milisegundo de tiempo $t=0,001seg$ y ¿Cuál es el valor de la fem de auto-inducción después de 1 minuto?
    \item Con los datos del problema anterior elaborar un gráfico donde en el eje vertical(ordenadas) se coloque el valor de f.e.m auto-inducción $e_{auto}$  en el eje de las abscisas los valores de tiempo para $t_1=0.01s$, $t_2=0.1s$, $t_3=0.5s$, $t_4=0.8s$, $t_5=1s$, $t_6=2s$, $t_7=4s,t_8=5s$ y $t_9=10s$
    \item Una bobina de $N=200$ espiras se mueve cortando perpendicularmente un campo magnético. La variación de flujo experimentado es uniforme y va desde los $\Phi_1=5mWb$ a los $\Phi_2=10mWb$ en un intervalo de tiempo de $\Delta t = 0.5 s$.Averiguar la f.e.m inducida.
    \item Una bobina de $N=100$ espiras corta con un ángulo de $45^o$ un campo magnético uniforme, la variación de flujo experimentada es uniforme y es de $\Delta \Phi = 4mWb$ en $\Delta t = 1s$. Averiguar la f.e.m inducida.
    \item Los terminales de una bobina son conectados a una resistencia $R=1\Omega$ y se genera una corriente de $I=0,5A$ y cuando esta se mueve produce una variación de flujo de  $\Delta \Phi =5mWb$ en $\Delta t = 2s$ ¿Cuántas espiras N tiene?
    \item Hallar la longitud que tiene que tener un conductor que corta perpendicularmente un campo magnético el cual posee una inducción de B=3T, y moviéndose a una velocidad de $v=5m/s$ genera una f.e.m de $e=10V$.
    \item Calcular el valor de la f.e.m. de autoinducción que desarrollará una bobina con un coeficiente de autoinducción de $L=50mH$, si se aplica una corriente que crece regularmente desde cero hasta $I=10A$ en un tiempo de $t=0.01s$.
    \item Calcular el coeficiente de autoinducción si se genera una f.e.m de $e=60V$ con una corriente que va desde $I_i=0A$ a $I_f=12A$ en el tiempo de $\Delta t=0,1s$.
    \item Calcular cuánto debe variar la corriente de forma uniforme para que se produzca una f.e.m. de $e=12V$ en una bobina con coeficiente de autoinducción de $L=60mH$, en el tiempo de $\Delta t=0,1s$.
    \item Una bobina posee $N=500$ espiras y produce un flujo magnético de $\Phi=10mWb$ cuando es atravesada por una corriente de $I=10A$. Determinar el coeficiente de autoinducción.
    \item Una bobina que posee $N=500$ espiras tiene un coeficiente de autoinducción $L=0,5H$. Averiguar si produce un flujo magnético de $\Phi=10mWb$ ¿Cuánta corriente circula por la misma?
    \item Averiguar el número de espiras que tiene una bobina que produce una variación de flujo magnético de $\Delta \Phi=0,01Wb$ cuando es travesada por una corriente de 10A y tiene un coeficiente de autoinducción igual a $L=0,5H$.
    \item Un conductor de $l=25cm$ se mueve en forma circular dentro de un campo magnético cortando el flujo magnético según la función $sen(\alpha)$ si el flujo máximo $\Phi=50mWb$. Si demora $t=1s$ en dar una vuelta. Realizar un gráfico donde el eje de las ordenadas sea el tiempo en segundos y las ordenadas el valor de la f.e.m generada.
    \item Un imán produce una inducción $\Phi=100mWb$ en una superficie de $s=20cm^2$ determinar el valor de la inducción del imán.
    \item Un super imán de $B=10T$ tiene sección hexagonal de cuyo radio que lo circunscribe es $r=10cm$. Determinar el flujo magnético en una cara del imán $\Phi$.
    \item Un imán de $B=3T$ tiene una longitud de 50cm, se pasa sobre el imán un conductor cuyo longitud es $l=20cm$ con una velocidad del $v=5m/s$. Determinar el valor de la f.e.m. que se produce.¿A que velocidad hay que hacerlo pasar para obtener 10V?
    \item Diseñar un generador para que genere una f.em. 10V, cuando posea un conductor que mueva a la velocidad de 3.14m/s.
    
 \end{enumerate}
\end{document}
